\section{Introduction}
\label{sec:intro}

Long-term time-series forecasting (LTSF) underpins decision-making in energy,
weather, traffic, and industrial monitoring, where models must predict hundreds
of future steps from a limited lookback. The dominant paradigm for several years
was the Transformer, with a succession of increasingly elaborate architectures
\citep{zhou2021informer,wu2021autoformer,zhou2022fedformer,nie2023patchtst}.
This trend was sharply questioned by the LTSF-Linear family
\citep{zeng2023dlinear}, which showed that embarrassingly simple linear models
match or exceed elaborate Transformers on the standard benchmarks, and by the
subsequent analysis of RLinear \citep{li2023rlinear}, which attributed much of
the linear models' strength to two ingredients: channel independence and
reversible instance normalization (RevIN) \citep{kim2022revin}. These findings
reframed LTSF as a setting where a careful, lightweight model can be both more
accurate and orders of magnitude cheaper than a heavy one---a goal directly
aligned with the Green~AI agenda \citep{schwartz2020greenai,strubell2019energy},
which argues that parameter, compute, and energy cost should be first-class,
reported evaluation criteria.

Despite their strength, linear LTSF models leave two specific gaps that motivate
this work. \textbf{First}, RevIN---now a near-universal plug-in---de-normalizes
the \emph{entire} forecast horizon with the \emph{single} mean and standard
deviation computed over the lookback. Under non-stationarity, the appropriate
level and scale drift across the horizon: the lookback statistics are a good
estimate for the first predicted step and progressively worse for the last.
RevIN's fixed, horizon-agnostic inverse cannot express this drift, and the
resulting error grows with horizon distance. \textbf{Second}, the decomposition
used by DLinear---a \emph{fixed} moving-average kernel that splits the series
into trend and remainder---is a single, hand-set low-pass filter applied in the
time domain, with no ability to adapt its band structure to a dataset's spectral
content, and it forecasts the remainder with an unconstrained head rather than
treating high-frequency structure as a first-class, separately modeled
component.

We address both gaps with \method, an ultra-lightweight, channel-independent
\emph{frequency-decomposed linear} forecaster. The backbone replaces the single
linear head with $K$ per-band heads fed by a learnable, lossless,
partition-of-unity spectral decomposition that adds only $2(K{-}1)$ scalar
parameters and generalizes DLinear's fixed moving-average split into a learnable
frequency-domain split. We are careful not to overclaim: we do \emph{not} assert
\method\ is the first frequency-domain linear forecaster. The closest prior art,
FITS \citep{xu2024fits}, also combines RevIN with a frequency-domain linear map,
but it applies a low-pass cutoff and \emph{discards} the high-frequency content
above it; \method\ instead \emph{keeps} all frequency content and routes the
high-frequency band to its own dedicated head, so residual structure is
\emph{modeled} rather than thrown away. Empirically, \method\ is the best
lightweight model on the standard long-term forecasting benchmarks, and---most
notably---at long lookback ($L{=}336$) it attains a lower average error than a
PatchTST Transformer \citep{nie2023patchtst} while using about $4\times$ fewer
parameters, $2.2\times$ less peak GPU memory, and $2.2\times$ less time per
epoch, all on a single 4\,GB laptop GPU.

On top of this backbone we introduce \emph{Adaptive Reversible Instance
Normalization} (\arevin), a \emph{regime-adaptive} reversible normalization that
strictly generalizes RevIN. Rather than re-injecting one fixed statistic at every
step, \arevin\ applies a small learnable per-step scale and shift plus a term that
extrapolates the \emph{observed} lookback drift across the horizon, all behind a
learnable gate; RevIN is recovered \emph{exactly} when the gate is closed.
Empirically, \arevin's benefit is regime-dependent, exactly as its name suggests:
it engages under non-stationarity---reducing error by up to about $5\%$ on the
strongly non-stationary ILI dataset, monotonically in the learned gate---while on
stationary benchmarks it reduces to RevIN with no harm. Our analysis also yields a
methodological insight: because \arevin's adaptive parameters start at their
identity values, the gate is gradient-starved at the conventional small
initialization and stays dormant even where adapting would help; initializing the
gate at the identity point (so its gradient is non-zero) lets it engage, while
still recovering RevIN exactly at the closed-gate value. \arevin\ is
philosophically aligned with the de-stationarization idea of Non-stationary
Transformers \citep{liu2022nonstationary}---removed statistics should be
re-injected \emph{adaptively}---but is realized as an attention-free module suited
to a linear backbone.

Because each component is an opt-in generalization that degenerates to a known
baseline, \method\ admits the clean theoretical framing
$\text{Linear} \subseteq \text{RLinear} \subseteq \method$, with DLinear's
decomposition recovered as a fixed-mask special case; this strict nesting makes
every added component independently ablatable. The entire study---training,
ablations, and a dedicated non-stationarity analysis---is run on a single 4\,GB
laptop GPU and is fully reproducible, with all baselines re-implemented and run by
us under identical splits, normalization, seeds, and hardware, and with
parameters, FLOPs, training time, and peak memory reported alongside MSE and MAE.
Consistent with our honesty bar, we report that the accuracy gains over RLinear on
the standard (stationary) benchmarks are modest, and we are explicit that the
headline strengths there are efficiency and the long-lookback win over the
Transformer.

\paragraph{Contributions.} In summary, this paper makes three contributions:
\begin{itemize}
  \item \textbf{\method.} An ultra-lightweight frequency-decomposed linear
        forecaster---a learnable, lossless band split feeding per-band linear
        heads---that is the best lightweight model on standard LTSF benchmarks
        and, at long lookback ($L{=}336$), surpasses the PatchTST Transformer at
        roughly $4\times$ fewer parameters and $2.2\times$ less memory on a single
        4\,GB laptop GPU, with improvements that are statistically significant
        under paired Wilcoxon tests ($p \approx 10^{-6}$).
  \item \textbf{A-RevIN.} A regime-adaptive reversible normalization that
        generalizes RevIN, together with an analysis showing (i) it must be
        initialized to avoid a gradient trap, and (ii) it engages and helps under
        non-stationarity---validated on both a real strongly non-stationary
        dataset (ILI) and a controlled synthetic drift sweep---while degrading
        gracefully to RevIN on stationary data.
  \item \textbf{A reproducible accuracy--efficiency study.} A fully reproducible
        study on commodity hardware (Green~AI), with ablations isolating each
        component and a dedicated non-stationarity analysis.
\end{itemize}
